\documentclass[11pt,a4paper]{article}
\usepackage[T1]{fontenc}
\usepackage[utf8]{inputenc}
\usepackage[scaled=0.95]{helvet}
\renewcommand{\familydefault}{\sfdefault}
\usepackage[top=1.7cm,bottom=1.5cm,left=2.2cm,right=2.2cm]{geometry}
\usepackage{microtype}
\usepackage[hidelinks]{hyperref}
\pagestyle{empty}
\setlength{\parindent}{0pt}
\setlength{\parskip}{7pt}

\begin{document}

{\large\bfseries Grover Aruquipa}\par\vspace{-6pt}
{\small Qu\'ebec, QC, Canada \ \textbullet\ \ grover.aruquipa.9@gmail.com \ \textbullet\ \ groveraruquipa.github.io}

\vspace{4pt}
R\&D Team, Distalmotion SA, \'Epalinges\\
R\&D Mechanical Engineer

Dear Distalmotion R\&D Team,

For the last three years I have been working on one mechanical problem: how to roll a tool without limit while still opening and closing a jaw at its tip, with every motor kept at the base. A wristed surgical instrument faces the same problem, only smaller and with a patient on the other side. When I read that you are developing the next generation of Dexter, I wanted to apply right away.

I am finishing my PhD at Universit\'e Laval with Prof.\ Cl\'ement Gosselin and Prof.\ Patrice Lambert, and I defend in early 2027. I have designed and built three robots with 4, 7 and 8 actuators. For each one I did the whole cycle myself: the requirements and the concept, the SolidWorks design of the joints and transmissions, the choice of motors, encoders, gearboxes and bearings, the FEA of the critical parts, the machining and printing, the assembly and wiring, and the test. When a part did not work, I took it apart on the bench, found out why and redesigned it.

The robot I would most like to show you is a 4-DoF coaxial wrist-gripper. It sends four actuation channels through concentric hollow shafts in a 180\,mm footprint, with 0.25\,kg of moving mass and a non-backdrivable jaw. On the bench it turned its tool ten full revolutions with less than one degree of tracking error while holding the object, under 500\,Hz real-time control. Right now I am building a Delta robot with an integrated Sarrus gripper that estimates the grasp force from the motor currents, without a sensor at the jaw, so it can hold fragile objects without crushing them. Controlling force at the tip without putting a sensor there is exactly the kind of problem I want to solve for Dexter.

Your posting asks for someone who can lead a mechanical subsystem from the early concept to a verified product, and that is exactly how I have worked for the last three years. I write the requirements before the CAD, I design every part for the process that will make it, I build the setups that measure whether the mechanism meets its targets, and I do not consider a design finished until the test data say so. Every one of my designs has gone through review by my supervisors and by peer reviewers in the ASME Journal of Mechanisms and Robotics and Mechanism and Machine Theory, so I am used to defending each choice with numbers and keeping the documentation that supports it. I also lead two graduate students day to day on an 8-actuator robot, so I know how to split a mechanical system into subsystems, assign them, and bring them back together on the bench.

I live and work in French and English in Qu\'ebec, so working in the Lausanne area would feel natural, and I am ready to relocate and work on site in \'Epalinges from January 2027. Prof.\ Gosselin, Prof.\ Lambert and Univ.-Prof.\ Andreas M\"uller of JKU Linz, who hosted my research stay this year, can speak to my work; their details are in my CV.

Thank you for your time. I would be happy to show you the prototypes on video, or to bring the wrist-gripper to an interview.

Sincerely,\\
Grover Aruquipa

\end{document}
